\documentclass{article}
\usepackage[utf8]{inputenc}
\usepackage[T1]{fontenc}
\usepackage{geometry}
\usepackage{graphicx}
\usepackage{float}
\usepackage{hyperref}
\usepackage[english,indonesian]{babel}

\title{Optimasi Pemrosesan Template LaTeX Menggunakan Pendekatan Berbasis AI}
\author{Academic Support Agent}
\date{\today}

\begin{document}
\maketitle

\begin{abstract}

\end{abstract}

\chapter{PENDAHULUAN}


\chapter{TINJAUAN PUSTAKA}


% Original content (commented out):
% \documentclass{article} 
\usepackage{pgffor}
\usepackage{enumitem} 

\usepackage{multicol}
\usepackage{textgreek}
\makeatletter
\newcommand{\greekalpha}[1]{\c@greekalpha{#1}}
\newcommand{\c@greekalpha}[1]{%
  {%
    \ifcase\number\value{#1} %
    \or
    \textalpha
    \or
    \textbeta
    \or
    \textgamma
    \or
    \textdelta
    \or
    \textepsilon
    \or
    \textzeta
    \or
    \texteta
    \or
    \texttheta % or \straighttheta
    \or
    \textiota
    \or
    \textkappa
    \or
    \textlambda
    \or
    \textmu
    \or
    \textnu
    \or
    \textxi
    \or
    \textomikron
    \or
    \textrho
    \or
    \textpi
    \or
    \textsigma
    \or
    \texttau
    \or
    \textupsilon
    \or
    \textphi
    \or
    \textchi
    \or
    \textpsi
    \or
    \textomega
    \fi
  }%
}


\newcommand{\Greekalpha}[1]{\c@Greekalpha{#1}}
\newcommand{\c@Greekalpha}[1]{%
  {%
    \ifcase\number\value{#1} %
    \or
    \textAlpha
    \or
    \textBeta
    \or
    \textGamma
    \or
    \textDelta
    \or
    \textEpsilon
    \or
    \textZeta
    \or
    \textEta
    \or
    \textTheta
    \or
    \textIota
    \or
    \textKappa
    \or
    \textLambda
    \or
    \textMu
    \or
    \textNu
    \or
    \textXi
    \or
    \textOmikron
    \or
    \textRho
    \or
    \textPi
    \or
    \textSigma
    \or
    \textTau
    \or
    \textUpsilon
    \or
    \textPhi
    \or
    \textChi
    \or
    \textPsi
    \or
    \textOmega
    \fi
  }%
}



\AddEnumerateCounter*{\greekalpha}{\c@greekalpha}{5}

\AddEnumerateCounter*{\Greekalpha}{\c@Greekalpha}{5}
\makeatother
\usepackage[utf8]{inputenc}
\usepackage{amsmath}
\usepackage{enumitem}
\usepackage{systeme}
\usepackage{amssymb}
\usepackage{tikz} 
\usepackage{amsmath,amssymb,amsthm,amssymb}
\usepackage{arydshln}
\newcommand{\R}{\mathbb{R}}
\newcommand{\Pp}{\mathbb{P}}
\usepackage{datetime}
\newdateformat{mydate}{\shortmonthname[\the\month]  \the\year}

\usepackage{graphicx}

\makeatletter
\newcommand*\bigcdot{\mathpalette\bigcdot@{.5}}
\newcommand*\bigcdot@[2]{\mathbin{\vcenter{\hbox{\scalebox{#2}{$\m@th#1\bullet$}}}}}
\makeatother


\title{Differential Equations HW Chapter.Section}
\author{First Initial. Middle Initial. Surname }
\date{\today}
\begin{document}

\maketitle


%I use the Sator Square as my seal, please use something else. But this does work to show how the \array command works. 

\[\begin{array}{ccccc}
\text{R} & \text{O} & \text{T} & \text{A} & \text{S}\\
\text{O} & \text{P} & \text{E} & \text{R} & \text{A}\\
\text{T} & \text{E} & \text{N} & \text{E} & \text{T}\\
\text{A} & \text{R} & \text{E} & \text{P} & \text{O}\\
\text{S} & \text{A} & \text{T} & \text{O} & \text{R}  \\
\end{array}
\]



\newpage


\newpage
% You "section counter" should be set to one behind whatever problem you are on, such that the \section command gives what problem in the chapter you are working on. 
\setcounter{section}{0}
\section{Chapter.Section.Problem}
\subsection{Goal}
\subsection{Given}
%%%%%%%%%%%%%%%%%%%%%%%%%%%%%%%%%
%   below is an example of how to write a differential eq


$$\frac{\delta y}{\delta x}= P(x)y=Q(x) $$

%%%%%%%%%%%%%%%%%%%%%%%%%%%%%%%%%

\subsection{Therefore}
\textbf{\textit{Q.E.D}}


%if you have multiple sections in your p-set then use the below command to separate them again sticking to the Chapter.Section format 

\newpage

\begin{center}
\begin{frame}{}
  \centering \Huge
  \emph{Chapter.Section}
\end{frame}    
\end{center}

\newpage



\newpage
\setcounter{section}{0}
\section{Chapter.Section.Problem}
\subsection{Goal}
\subsection{Given}
\subsection{Therefore}
\textbf{\textit{Q.E.D}}



\textbf{\textit{Q.E.D}}


\end{document}



\end{document}